\section{System Design}
\label{sec:system}

\todo{architecture figure + longer prose.}

\sys{} sits between a LangGraph-orchestrated agent pool and a shared
ChromaDB vector memory. It exposes three control-plane components:

\paragraph{Behavioral Monitor (\cref{sec:methodology:signals}).}
For every PR-handling step, each agent emits six scalar signals:
tool-call entropy, decision latency, memory-query entropy, KL
divergence against a pre-attack baseline distribution, retrieval-hit
rate, and action-repetition rate. The monitor persists these signals
to SQLite (dev) or PostgreSQL (prod) for offline analysis and streams
them to the Temporal Health Scorer.

\paragraph{Temporal Health Scorer (\cref{sec:methodology:detector}).}
The scorer combines two detectors on the dominant signal
(KL divergence): a Bayesian Online Change-Point Detector that flags
regime shifts with calibrated run-length posteriors, and a Chronos-2
probabilistic forecaster that supplies an expected-range prior. A
positive detection from either branch decrements the agent's health
score; a sustained normal regime lets the score recover.

\paragraph{Decentralized Task Allocator (\cref{sec:methodology:alloc}).}
New PRs are offered via a lightweight negotiation protocol; each agent
bids a cost that is a function of its own health score and queue
depth. The allocator picks the lowest-cost bid. A drifting agent's
health score drops, its cost rises, and work routes away from it
automatically, without a centralized kill-switch.

\paragraph{Human Escalation Layer.}
When the aggregate health score across the pool drops below a
configurable threshold, or when change-points cluster in a short
window, the system raises a reviewer alert through the FastAPI
dashboard (Phase 8) rather than silently quarantining agents.
